% Generated draft: requires mathematical and source review.
\documentclass[11pt]{article}
\usepackage[margin=25mm]{geometry}
\usepackage{fontspec}
\setmainfont{Latin Modern Roman}
\usepackage{amsmath,amssymb,mathtools,unicode-math}
\setmathfont{Latin Modern Math}
\usepackage{enumitem,needspace,xurl,hyperref}
\setlength{\parindent}{0pt}
\setlength{\parskip}{5pt}
\setlength{\emergencystretch}{3em}
\newcommand{\N}{\mathbb{N}}
\newcommand{\Z}{\mathbb{Z}}
\newcommand{\Q}{\mathbb{Q}}
\newcommand{\R}{\mathbb{R}}
\newcommand{\C}{\mathbb{C}}
\newcommand{\F}{\mathbb{F}}
\newcommand{\A}{\mathbb{A}}
\newcommand{\PP}{\mathbb{P}}
\newcommand{\E}{\mathbb{E}}
\newcommand{\eps}{\varepsilon}
\newcommand{\dd}{\,\mathrm{d}}
\begin{document}
\section*{UnsolvedMath ID: 30002097}
\textbf{Source problem code:} catalogue entry
\subsection{Definitions and mathematical statement}
For terminology only, let \(\mathbb K\in\{\R,\C\}\), let \(m,n\geq1\) be integers, and let \(A\in\mathbb K^{m\times n}\). These choices are not specified for the scattering model in the supplied record. For an integer \(1\leq s\leq n\), define
\[
\Sigma_s(\mathbb K^n)=\{x\in\mathbb K^n:\#\{j:x_j\neq0\}\leq s\}.
\]
The matrix \(A\) satisfies the restricted isometry property of order \(s\) with constant \(\delta\in[0,1)\) if
\[
(1-\delta)\|x\|_2^2\leq\|Ax\|_2^2\leq(1+\delta)\|x\|_2^2
\qquad\text{for every }x\in\Sigma_s(\mathbb K^n).
\]
For a fixed matrix \(A\), exact noiseless uniform recovery at sparsity \(s\) means that there exists a single reconstruction map \(\Delta_A:\mathbb K^m\to\mathbb K^n\) such that
\[
\Delta_A(Ax)=x\qquad\text{for every }x\in\Sigma_s(\mathbb K^n).
\]
This definition explains one conventional meaning of uniform recovery; the source does not specify whether exact noiseless recovery, stable recovery, or recovery by a particular algorithm is intended. For a random matrix, uniformity requires the assertion to hold simultaneously for all relevant sparse signals on the same realization of the matrix; the required probability of such an event is not supplied.

\textbf{Supplied open question.} Does the random scattering matrix used for sparse recovery satisfy a restricted isometry property in a parameter regime strong enough to yield uniform recovery of all sparse signals?

\textbf{Unresolved formulation.} The supplied record does not define the scattering matrix, its dimensions, entry formula, probability law, normalization, underlying physical or geometric parameters, or admissible parameter regime. It also does not specify the signal field, sparsity range, required restricted isometry order and constant, recovery method, noise assumptions, or probability bound. Consequently, the question above cannot be converted from this record into a uniquely quantified conjecture about a specified random matrix ensemble. The definitions describe the terminology without selecting the missing model or recovery target. The supplied designation as open is a source label, not an independently verified mathematical conclusion.
\subsection{Short English statement}
The question asks whether the intended random scattering matrices preserve the Euclidean lengths of sparse vectors sufficiently well to support recovery simultaneously for every signal of the prescribed sparsity. The record does not specify the matrix ensemble or the quantitative recovery target, so a precise conjecture remains unresolved and requires human mathematical review.
\subsection{Sources}
\begin{itemize}
\item Supplied local UnsolvedMath registry record, permanent ID \(30002097\), source problem code \texttt{catalogue entry}, title \emph{Restricted Isometries for Random Scattering Matrices}; collection: Oberwolfach Reports Open Problems. The project identifies \texttt{data/unsolvedmath/problems.json} as the primary registry and \texttt{catalog/problems.json} as the additional-record catalog. This draft uses only the supplied complete record.
\item Supplied report citation: \emph{Applied Harmonic Analysis and Sparse Approximation (2013)}, Oberwolfach Reports, DOI 10.4171/owr/2012/29, pp.\ 1759--1843. Original OWR report: \url{https://doi.org/10.4171/owr/2012/29}. The background attributes the question to collected workshop open-problem discussions. No individual proposer is supplied. The proposed year and year tag are 2012, while the supplied report title contains 2013; this dating discrepancy is unresolved.
\item Supplied literature source labeled ``modified CoSaMP electromagnetic imaging'': \url{https://repository.kaust.edu.sa/items/6801b3b7-0d36-44f0-84fd-cdeee9e0a8c3}. No authors, publication date, or full bibliographic citation for this item are supplied.
\item Literature assessment dated exactly \texttt{2026-08-21}: the supplied curation labels the problem open and states that nonuniform sparse-recovery results and numerical damping methods are available, but that the RIP estimate for the specific scattering matrix remains unestablished and the matrix model is not generalized enough for a theorem covering the claim. This is supplied dated curation derived from the OWR material, not an independently verified assessment. The record itself requests independent verification of time-sensitive status and cited claims.
\item The record labels its statement ``exact'' and reports manual verification that the extracted statement is self-contained, mathematically coherent, and free of LaTeX or report-layout artifacts. These are supplied verification claims; they do not resolve the missing matrix model and quantitative hypotheses identified here. No source was independently visited or verified for this draft.
\end{itemize}
\end{document}
